% Pursuing Stacks (cotes 134-2 to 134-8): archive cote 134-6, site page 22 (PDF page 23), typescript page 341.
% typescript. Draft compared: ps5.tex:2466-2521.
% AI-transcribed by Claude Opus 5.5 (claude-opus-5-5, adaptive thinking, effort high), Message Batches API, batch ps3 (msgbatch_01XdFmuVVBAcgc7EFxdtCqsj), 2026-11-12; not yet checked by eye.
\scanpage{134-6/22}
\mspage{341}
arbitrary semisimplicial sets,  which took me a few weeks to work out,
and \inserted{alongside} getting back into homotopy and cohomology formalism again.  I had
the feeling that this structure,  or the technically more adequate
dual \xout{xx} ``coalgebra'' structure,  might well turn out to be \unsure{the} more re-%
fined version of (27)  needed for recovering homotopy types - \xout{xx}
or at any rate 1-connected ones.  I gave a talk about the matter at
IHES while things were still hot in my mind  -  but it does'nt seem
it went really through.  \xout{Maybe} It does'nt seem this structure (which was
worked out independently by someone else too,  I understand) \deleted{got} \inserted{has become} a fami-%
liar notion to topologists.  Maybe one reason is that \inserted{most} topologists \xout{never}
and homotopy theorists never really got acquainted with the formalism
of derived categories  -  and it seems that \inserted{moreover,} by the mid-sevent\unsure{ie}s,  it
had \inserted{even,} become altogether unfashionable and ``mal vu'' to make any mention of
them let alone work with them,  \xout{even} also \deleted{for} \inserted{among} some of the people who
\inserted{during some time} had been helping \xout{during some time} develop it.  \unsure{NOw|Now} one of the main
points I was making in that talk was a somewhat delicate property of
derived categories of abelian categories,  with respect to binomial
coefficients  -  too bad !

I have not heard since about any work done in th\unsure{is} direction I
am reflecting about now (somewhat retrospectively)  -  namely recovering
homotopy types from their abelianization,  \uline{plus} extra structure.  For
all I know,  \uline{the}  relevant structure may well be the differential alge-%
bra \xout{structure} with divided powers  structure embodied by the De Rham
complex (with a bigraduation however instead of just a graduation),  or
it's coalgebra version  -  viewed as defining an object of a suitable
derived category.  (Of course,  when there is a fundamental group \hand{$\pi$}
around,  one will have to look at a slightly more complex structure
still,  involving operations of \hand{$\pi$},  by looking at the De Rham complex
of the universal covering.)   If it is just the matter of describing
homotopy types in terms of \xout{xxxxx} other models than semisimplicial com-%
plexes,  it must be admitted that the new models are \inserted{of} uncomparably more
intricate description than the complexes\over{.}{!}  There \uline{is} however one feature
of it which greatly struck me by that time,  and still seems to me quite
intriguing,  namely \uline{that this structure},  \uline{although definitely not}
\uline{``abelian'' any more}  (due to multiplication as well as to divided power
structure),  \uline{makes a sense over any commutative ground ring} (or even
scheme etc).  When this ring is $\underline{Q}$ ,  the ``models'' we get modelize
rational homotopy types,  which was the starting point of my reflecti-%

\margin{22}
